\documentclass[a4paper,12pt]{article}
\usepackage[T2A]{fontenc}
\usepackage[utf8]{inputenc}
\usepackage[russian]{babel}
\usepackage{amsmath,amssymb,mathtools}
\usepackage{PTSans}
\renewcommand{\familydefault}{\sfdefault}
\usepackage{icomma}
\usepackage{geometry}
\usepackage{microtype}
\usepackage{tikz}
\usetikzlibrary{calc,arrows.meta,positioning,shapes.geometric}
\usepackage{tabularx,booktabs}
\usepackage{enumitem}
\usepackage{parskip}
\usepackage{fancyhdr}
\usepackage{setspace}
\usepackage{needspace}
\usepackage{xcolor}

\geometry{left=20mm,right=20mm,top=19mm,bottom=20mm}
\setlength{\headheight}{25pt}
\onehalfspacing
\setlength{\parindent}{0pt}
\setlist[itemize]{leftmargin=7mm,itemsep=1pt,topsep=3pt}
\setlist[enumerate]{leftmargin=8mm,itemsep=4pt,topsep=4pt}
\definecolor{accent}{HTML}{147D7E}
\definecolor{darkaccent}{HTML}{0B5052}
\definecolor{textgray}{HTML}{353C3D}
\definecolor{warn}{HTML}{8B5946}
\pagestyle{fancy}
\fancyhf{}
\fancyhead[L]{\small\color{textgray}НОД, НОК и алгоритм Евклида}
\fancyhead[R]{\small\color{textgray}08.10.26}
\fancyfoot[C]{\small\color{textgray}\thepage}
\renewcommand{\headrulewidth}{0.3pt}
\renewcommand{\headrule}{\hbox to\headwidth{\color{accent}\leaders\hrule height \headrulewidth\hfill}}
\newcommand{\topic}[1]{%
 \Needspace{5\baselineskip}%
 \vspace{0.6em}{\Large\bfseries\color{darkaccent}#1}\par
 \vspace{0.15em}{\color{accent}\rule{\linewidth}{0.55pt}}\par\vspace{0.2em}}
\newcommand{\key}[1]{\textbf{\color{darkaccent}#1}}
\newcommand{\warning}[1]{\par\textbf{\color{warn}Внимание.} #1\par}
\newcommand{\checkitem}{\(\square\)}
\renewcommand{\today}{08.10.26}
\begin{document}
\thispagestyle{empty}
\vspace*{23mm}
\begin{center}
 {\fontsize{31}{35}\selectfont\bfseries\color{darkaccent}Чек-лист}\\[5mm]
 {\fontsize{21}{26}\selectfont\bfseries НОД и НОК}\\[4mm]
 {\Large Алгоритм Евклида и первые олимпиадные рассуждения}\\[15mm]
 {\large Ксения Клыкова}\\[3mm]
 {\large \today}\\[17mm]
 \begin{minipage}{0.8\textwidth}\centering\color{textgray}
 Научиться не только вычислять НОД и НОК, но и замечать
 ограничения на числа, доказывать невозможность и находить
 подходящий пример.
 \end{minipage}
\end{center}
\vfill
\begin{center}
 {\large\bfseries\color{darkaccent}Лёвин Артём Александрович}\\[2mm]
 {\large\color{textgray}эксклюзивно для Ксении Клыковой}
\end{center}
\vspace{8mm}
\begin{tikzpicture}[remember picture,overlay]
 \draw[accent,line width=1.2pt]
 ([xshift=18mm,yshift=17mm]current page.south west)
 .. controls ([xshift=65mm,yshift=28mm]current page.south west)
 and ([xshift=-62mm,yshift=7mm]current page.south east)
 .. ([xshift=-18mm,yshift=20mm]current page.south east);
\end{tikzpicture}
\clearpage

\topic{1. Основные понятия и разложение на простые множители}
\key{Наибольший общий делитель (НОД)} двух положительных целых чисел --- наибольшее положительное число, на которое делится каждое из них.

\key{Наименьшее общее кратное (НОК)} --- наименьшее положительное число, которое делится на каждое из них.

\begin{itemize}
 \item \checkitem{} Разложите числа на простые множители.
 \item \checkitem{} Для \key{НОД} возьмите каждый общий простой множитель в \emph{меньшей} степени.
 \item \checkitem{} Для \key{НОК} возьмите каждый встречающийся простой множитель в \emph{большей} степени.
\end{itemize}

Пример:
\[
24=2^3\cdot3,\qquad32=2^5,
\]
\[
\boxed{\gcd(24,32)=2^3=8},\qquad
\boxed{\operatorname{lcm}(24,32)=2^5\cdot3=96}.
\]
Здесь и далее $\gcd(a,b)$ обозначает НОД, а $\operatorname{lcm}(a,b)$ --- НОК. В краткой записи также употребляют $(a,b)$ и $[a,b]$, когда смысл понятен из контекста.

\topic{2. Алгоритм Евклида: НОД через остатки}
Для положительных целых $a\ge b$ деление с остатком имеет вид
\[
a=qb+r,\qquad 0\le r<b,
\]
где $q$ --- неполное частное, $r$ --- остаток. Главное свойство:
\[
\boxed{\gcd(a,b)=\gcd(b,r)}.
\]
Повторяйте деление: новый делитель становится делимым, а ненулевой остаток --- новым делителем. \key{Последний ненулевой остаток равен НОД.}

\key{Пример.} Найдём $\gcd(308,112)$:
\[
\begin{aligned}
308&=2\cdot112+84,\\
112&=1\cdot84+28,\\
84&=3\cdot28+0.
\end{aligned}
\qquad\Longrightarrow\qquad \boxed{\gcd(308,112)=28}.
\]
\warning{Останавливайтесь при нулевом остатке. Сам ноль ответом не является.}

\clearpage
\vspace*{17mm}
\thispagestyle{fancy}
\begin{center}
 {\LARGE\bfseries\color{darkaccent}Алгоритм Евклида}\par\vspace{4mm}
 {\color{textgray}Последовательно заменяем пару чисел меньшим числом и остатком.}
\end{center}
\vspace{10mm}
\begin{center}
\begin{tikzpicture}[
 >=Stealth, line width=.85pt,
 term/.style={draw=darkaccent,ellipse,minimum width=46mm,minimum height=11mm,align=center,fill=white},
 action/.style={draw=darkaccent,rounded corners=2mm,minimum width=53mm,minimum height=15mm,text width=49mm,align=center,fill=white},
 test/.style={draw=darkaccent,diamond,aspect=2.3,inner sep=1mm,text width=28mm,align=center,fill=white},
 note/.style={font=\small\bfseries,fill=white,inner sep=1pt},
 arrow/.style={->,draw=accent,line width=1.1pt}
]
\node[term] (start) at (0,0) {Начало: $a,b>0$};
\node[action] (norm) at (0,-2.0) {Если нужно, поменяйте числа местами: $a\ge b$};
\node[action] (divide) at (0,-4.4) {Разделите с остатком: $a=qb+r$, $0\le r<b$};
\node[test] (test) at (0,-7.0) {$r=0$?};
\node[term] (end) at (0,-9.65) {Ответ: НОД $=b$};
\node[action] (again) at (6.0,-7.0) {Замените пару: $a\gets b$, $b\gets r$};
\draw[arrow] (start) -- (norm);
\draw[arrow] (norm) -- (divide);
\draw[arrow] (divide) -- (test);
\draw[arrow] (test) -- node[note,left]{да} (end);
\draw[arrow] (test) -- node[note,above]{нет} (again);
\draw[arrow] (again.north) |- ([yshift=2mm]divide.east) -- (divide.east);
\end{tikzpicture}
\end{center}
\clearpage

\topic{3. Как распознать НОД и НОК в текстовой задаче}
\key{Одинаковые наборы без остатка.} Если $84$ помидора и $112$ огурцов нужно распределить по наибольшему числу одинаковых упаковок, то
\[
\gcd(84,112)=28.
\]
Получатся $28$ упаковок, в каждой $84:28=3$ помидора и $112:28=4$ огурца.

\key{Первое одновременное повторение.} Если два действия повторяются каждые $12$ и $18$ минут и совпали в начальный момент, то следующее совпадение наступит через
\[
\operatorname{lcm}(12,18)=36\text{ минут}.
\]
\warning{Слова «упаковки» и «встретятся» сами по себе не гарантируют выбор метода. Проверьте, что требуются максимальное число равных групп или первое совместное кратное.}

\topic{4. Главная связь между НОД и НОК}
Для любых положительных целых чисел $a,b$:
\[
\boxed{\gcd(a,b)\cdot\operatorname{lcm}(a,b)=ab}.
\]
Поэтому, зная три величины, часто можно найти четвёртую без отдельного разложения на множители.

\key{Пример.} Для $a=84$, $b=180$ и $\gcd(84,180)=12$:
\[
\operatorname{lcm}(84,180)=\frac{84\cdot180}{12}=1260.
\]
Почему формула работает? Для каждой простой основы $p$ показатели степеней удовлетворяют
\[
\min(\alpha,\beta)+\max(\alpha,\beta)=\alpha+\beta.
\]
\begin{itemize}
 \item \checkitem{} В формуле НОД и НОК относятся к \emph{одной и той же паре} $a,b$.
 \item \checkitem{} Следите, чтобы все найденные числа оставались положительными целыми.
\end{itemize}

\topic{5. Можно ли задать НОД и НОК произвольно?}
\key{Необходимое условие:} НОД любых двух положительных целых чисел всегда делит их НОК:
\[
\boxed{\gcd(a,b)\mid\operatorname{lcm}(a,b)}.
\]
Например, условие $\gcd(a,b)=14$ и $\operatorname{lcm}(a,b)=90$ \key{невыполнимо}, так как $90$ не делится на $14$.

И наоборот: если положительное $d$ делит положительное $m$, то пара $a=d$, $b=m$ даёт НОД $d$ и НОК $m$. Значит, это условие \emph{достаточно} для существования хотя бы одной пары.

\key{Приём «оценка + пример»:} сначала доказать ограничение для всех возможных чисел, затем предъявить конкретные числа, которые этому ограничению удовлетворяют. Если ограничение нарушено, достаточно доказательства невозможности.

\topic{6. Первые рассуждения олимпиадного типа}
\key{Задача с неизвестным числом.} Пусть
\[
n=2^x3^y5^z7^w,\quad x,y,z,w\in\mathbb Z_{\ge0},
\]
причём $\gcd(n,360)=120$, $\operatorname{lcm}(n,360)=2520$. Тогда
\[
 n=\frac{120\cdot2520}{360}=\boxed{840}
 =2^3\cdot3\cdot5\cdot7.
\]
Следовательно, $x=3$, $y=z=w=1$. Проверка исходных условий подтверждает результат.

\key{Задача с простым числом.} Найти простое $p$, если $\operatorname{lcm}(18,p)=90$.
Поскольку $p$ --- простой делитель $90=2\cdot3^2\cdot5$, возможны лишь $p=2,3,5$. При $p=2$ или $3$ НОК равен $18$; при $p=5$ --- $90$. Итак, $\boxed{p=5}$.

\key{Что важно заметить:}
\begin{itemize}
 \item \checkitem{} Сначала запишите формулу или ограничение и лишь затем вычисляйте.
 \item \checkitem{} Используйте дополнительные условия: «простое», «целое», «наименьшее», «все».
 \item \checkitem{} Для существования предъявите пример; для невозможности найдите противоречие.
 \item \checkitem{} Найденное из произведения число проверьте по исходным НОД и НОК.
\end{itemize}

\warning{Число $1$ простым не является. Простое число имеет ровно два различных положительных делителя.}


\clearpage
\topic{Тренировочный блок --- Путь Б}
\textbf{09--12 октября 2026 года. По 3 задания каждый день.} Все задания средней сложности. Записывайте ход рассуждений, особенно там, где требуется доказательство или перечисление всех вариантов. Ответы помещены после тренировки.

\vspace{2mm}
\topic{09.10.26 --- пятница}
\begin{enumerate}
 \item Найдите $\gcd(672,462)$ алгоритмом Евклида. Запишите все деления с остатком.
 \item Найдите $\operatorname{lcm}(84,126,180)$ через разложение на простые множители.
 \item Для натурального $n$ известно, что $\gcd(n,72)=18$ и $\operatorname{lcm}(n,72)=360$. Найдите $n$ и проверьте оба условия.
\end{enumerate}

\vspace{4mm}
\topic{10.10.26 --- суббота}
\begin{enumerate}
 \item Из $168$ карандашей, $210$ тетрадей и $294$ закладок составляют максимально возможное число одинаковых наборов без остатка. Сколько наборов получится? Сколько предметов каждого вида будет в одном наборе?
 \item Три сигнала подаются каждые $18$, $24$ и $40$ минут. В $08{:}00$ они прозвучали одновременно. Когда они впервые прозвучат одновременно снова? Сколько раз сработает каждый сигнал за прошедшее время (не считая начального момента)?
 \item Найдите \emph{все неупорядоченные пары} натуральных чисел $(a,b)$, для которых $\gcd(a,b)=18$ и $\operatorname{lcm}(a,b)=540$. Обоснуйте, что других пар нет.
\end{enumerate}

\clearpage
\topic{11.10.26 --- воскресенье}
\begin{enumerate}
 \item Выясните, какие данные о двух натуральных числах могут быть верными: \\
 (а) $\gcd(a,b)=18$, $\operatorname{lcm}(a,b)=450$; \\
 (б) $\gcd(a,b)=18$, $\operatorname{lcm}(a,b)=455$. \\
 Для возможного случая приведите конкретную пару чисел; невозможность обоснуйте.
 \item Число $n=2^x3^y5^z7^w$, где $x,y,z,w$ --- целые неотрицательные числа, удовлетворяет условиям $\gcd(n,420)=84$ и $\operatorname{lcm}(n,420)=1260$. Найдите $n$ и все четыре показателя степени. Проверьте результат.
 \item Найдите все простые числа $p$, для которых выполняется хотя бы одно из равенств:
 \[
 \operatorname{lcm}(84,p)=84\quad\text{или}\quad\operatorname{lcm}(84,p)=420.
 \]
 Укажите, для какого равенства подходит каждое найденное число.
\end{enumerate}

\vspace{3mm}
\topic{12.10.26 --- понедельник}
\begin{enumerate}
 \item Найдите \emph{все} натуральные числа $n$, для которых
 \[
 \gcd(n,60)=60,\qquad \operatorname{lcm}(n,36)=180.
 \]
 Объясните, почему список полный.
 \item Для натуральных $a,b$ известно: $\gcd(a,b)=12$, $\operatorname{lcm}(a,b)=420$. Найдите пару, у которой разность $|a-b|$ наименьшая. Обоснуйте выбор.
 \item Найдите наименьшее натуральное число $N$, которое кратно $24$ и $36$, а при делении на $5$ даёт остаток $3$. Покажите, почему найденное число действительно наименьшее.
\end{enumerate}

\clearpage
\topic{Ответы к тренировке --- Путь Б}
\renewcommand{\arraystretch}{1.32}
\begin{tabularx}{\linewidth}{@{}p{15mm}p{10mm}X@{}}
\toprule
Дата & № & Ответ и краткий ориентир \\
\midrule
\multicolumn{3}{@{}l}{\bfseries 09.10.26 --- пятница} \\
 & 1 & $42$: $672=462+210$, $462=2\cdot210+42$, $210=5\cdot42$. \\
 & 2 & $1260=2^2\cdot3^2\cdot5\cdot7$. \\
 & 3 & $n=90$; НОД $=18$, НОК $=360$. \\
\addlinespace[2mm]
\multicolumn{3}{@{}l}{\bfseries 10.10.26 --- суббота} \\
 & 1 & $42$ набора: по $4$ карандаша, $5$ тетрадей и $7$ закладок. \\
 & 2 & $14{:}00$; за $360$ минут сигналы сработают $20$, $15$ и $9$ раз. \\
 & 3 & $(18,540)$, $(36,270)$, $(54,180)$, $(90,108)$. Если $a=18u$, $b=18v$, то $uv=30$ и $\gcd(u,v)=1$. \\
\addlinespace[2mm]
\multicolumn{3}{@{}l}{\bfseries 11.10.26 --- воскресенье} \\
 & 1 & (а) возможно: $(18,450)$; (б) невозможно, так как $18\nmid455$. \\
 & 2 & $n=252=2^2\cdot3^2\cdot7$; $x=2$, $y=2$, $z=0$, $w=1$. \\
 & 3 & Для НОК $=84$: $p=2,3,7$; для НОК $=420$: $p=5$. \\
\addlinespace[2mm]
\multicolumn{3}{@{}l}{\bfseries 12.10.26 --- понедельник} \\
 & 1 & $n=60$ или $180$: из условий следует $60\mid n$ и $n\mid180$; оба числа подходят. \\
 & 2 & $(60,84)$; минимальная разность $24$. Возможные пары: $(12,420)$ и $(60,84)$. \\
 & 3 & $N=288$: число кратно $72=\operatorname{lcm}(24,36)$, а первые кратные $72,144,216,288$ дают остатки $2,4,1,3$ по модулю $5$. \\
\bottomrule
\end{tabularx}

\vfill
\begin{center}
{\small\color{textgray}Лёвин Артём Александрович эксклюзивно для Ксении Клыковой}
\end{center}
\end{document}